\documentclass[11pt,a4paper]{article}

\usepackage[utf8]{inputenc}
\usepackage[cyr]{aeguill}
\usepackage[francais]{babel}

% Les packages
%\usepackage[french]{babel}
\usepackage{fancybox}
\usepackage{graphics}
\usepackage{color}
\usepackage{latexsym}
\usepackage{amsmath,amsthm,amssymb}
\usepackage[plain]{fullpage} 
\usepackage{verbatim}
\usepackage{times,epsfig}
\usepackage{listings}
\usepackage{multirow}
\usepackage{amsmath,amsthm,amssymb,amscd,txfonts,bbm} % RAJOUT ICI !!!
\usepackage{graphics,epsfig} % RAJOUT ICI !!!
%\usepackage{algorithmic}
\usepackage[final]{pdfpages} 
\usepackage{fancyvrb}
\usepackage{enumerate}
\usepackage{listingsutf8}
\usepackage{comment}
\usepackage{verbatim}

% small et verbatim: to be coherent with \file
\newenvironment{smallverbatim}%
{\endgraf\small\verbatim}{\endverbatim}

\newenvironment{qv}
{\quote\Verbatim}
{\endquote\endVerbatim}

\usepackage{lastpage}

\newcommand{\terme}[1]{\texttt{#1}}
\newcommand{\termSet}[1]{\{\texttt{#1}\}}
\newcommand{\guill}[1]{<<\,#1\,>>}
% Macro pour l'inclusion de figure  
\newcommand{\fig}[2]{%
  \begin{figure*}[hbt]
   \begin{center}
      \includegraphics[width=12cm]{#1}
  \caption{\label{#1}  #2}
  \end{center}
 \end{figure*}
}
\newcommand{\figTex}[2]{%
  \begin{figure*}[hbt]
 \centerline{\input{#1.pstex_t}}
  \caption{\label{#1}  #2}
 \end{figure*}
}

% Macro résumés
\definecolor{turquoise}{rgb}{0.20 ,0.60, 0.60}
\definecolor{vert}{rgb}{0.40,  0.60 ,0.00}
\definecolor{violet}{rgb}{0.80 , 0.00 , 0.80}
\definecolor{bleu}{rgb}{0.20,  0.20,  0.80}
\definecolor{rose}{rgb}{1.00,  0.20,  0.40}

\lstset{basicstyle=\footnotesize,showstringspaces=false,language=XML,breaklines=false,columns=fullflexible,frame=shadowbox, captionpos=b}

\newenvironment{solution}[0]{
~\\ \color{red} \textbf{Solution : }\color{black}}{%
}
%\excludecomment{solution}


\newcommand{\cle}[1]{\textbf{#1}}
\def\fk#1{\textit{#1}}
\def\tble#1{\textit{#1}}

\parindent=0pt           
                             
\usepackage{url}
\usepackage{minted}

\newtheorem{Exercice}{Definition}


\begin{document}

\title{Sujet d'examen UE NFE204 - Session 1}
\author{Année universitaire 2025-2026}
\date{Examen session 1\,: 27 janvier 2026}

\maketitle


\begin{center}
{\bf \Large Consignes}\\
\begin{tabular}{| p{15cm} |}
\hline
Dur\'ee de l'examen : 3h00; Calculatrice autoris\'ee\,; Documents autorisés: 5 pages recto-verso de note personnelles.
Ce sujet contient \pageref{LastPage} pages.\\
\hline
\end{tabular}
\end{center}

\section*{Première partie : modélisation NoSQL (8 pts)}

Nous sommes en 2030, les attaques répétées contre
la presse indépendante ont fini par atteindre leur but\,: 
il n'y a plus de presse, de quotidiens ou de magazines. Il
ne reste que des articles engendrés par IA et 
quelques journalistes
indépendants, payés au caractère, 
qui écrivent des articles et les proposent à des plateformes de diffusion.
Ces journalistes sont (faiblement) rémunérés en fonction du nombre 
de personnes qui lisent
leur article sur ces plateformes.

Posons-nous quelques questions sur le fonctionnement
de cette presse ubérisée. Un modèle UML très sommaire est 
donné dans la figure\,\ref{exam-21-presseuber}. Les journalistes
publient des articles, et ces articles sont diffusés par des  plateformes (plusieurs plateformes 
si possible, qui elles-mêmes publient beaucoup d'articles).

\fig{presseuber}{Modèle (sommaire) de la diffusion d'articles ubérisés}

\vspace*{0.2cm}

Au départ, on représente tout cela avec une base relationnelle dont voici 
un tout petit échantillon.

\begin{center}

\begin{tabular}{l|p{6cm}|l}
\cle{ref} &  titre & année \\
\hline
ju21 & D. Trump propose de racheter la Corse & 2028 \\
ju798 & Les vaccins nuisent à la santé & 2030 \\
\hline
\end{tabular}

\centerline{La table \tble{Article}}
\end{center}

\begin{center}
\begin{minipage}[b]{8cm}
\begin{tabular}{l|l}
\cle{id} & nom \\
\hline
nt & N. Travers  \\
pr & P. Rigaux  \\
rfs & R. Fournier-S'niehotta \\
cdm & C. du Mouza \\
\hline
\end{tabular}

\centerline{La table \tble{Journaliste}}
\end{minipage} \hspace*{1cm} \
\begin{minipage}[b]{6cm}
\begin{tabular}{l|l}
refArticle  & idAuteur \\
\hline
ju21 & cdm \\
ju21 & pr \\
ju21  & nt \\
ju798 & pr \\
\hline
\end{tabular}
\centerline{La table des publications}
\end{minipage} 

\vspace*{0.2cm}
\begin{minipage}[b]{6cm}
\begin{tabular}{l|l}
\cle{idPlateforme} & nom\\
\hline
pj2 & Béru\\
pj5 & Croco \\
pj6 & Vautour \\
\hline
\end{tabular}
\centerline{La table des plateformes}
\end{minipage} \hspace*{1cm} \
\begin{minipage}[b]{6cm}
\begin{tabular}{l|l}s
refArticle  & idPlateforme \\
\hline
ju21 & pj2 \\
ju798 & pj6 \\
ju21  & pj6 \\
\hline
\end{tabular}
\centerline{La table des diffusions}
\end{minipage} 
\end{center}

\begin{enumerate}
  \item Proposez un format de document JSON représentant toutes les informations présentes dans les  
        tables ci-dessus et
        relatives à l'auteur \emph{P.~Rigaux} (représentation A).


\begin{solution}
    \begin{small}
    \begin{minted}{json}
      {"_id": "pr",
         "nom" : "P. Rigaux",
         "articles": [
           {"titre": "D. Trump propose de racheter la Corse", 
            "année": 2028
            "plateformes": ["Béru", "Vautour"]
            },
           {"titre": "Les vaccins nuisent à la santé", 
             "année": 2030
            "plateformes": ["Vautour"],
             }
           ]
        }
    \end{minted}
    \end{small}
\end{solution}

  \item Proposez un document JSON représentant toutes les 
        informations relatives à la plateforme \emph{Vautour} 
        présentes dans les tables ci-dessus
         (représentation B).
        

\begin{solution}
    \begin{small}
    \begin{minted}{json}
      {"_id": "pj6",
         "nom" : "Vautour",
         "articles": [
           {"titre": "D. Trump propose de racheter la Corse", 
            "année": 2028
            "auteurs": ["N. Travers", "P. Rigaux", "R. Fournier-S'niehotta"]
            },
           {"titre": "Les vaccins nuisent à la santé", 
             "année": 2030
            "auteurs": ["P.Rigaux"],
             }
           ]
        }
    \end{minted}
    \end{small}
\end{solution}


   \item Expliquez le calcul MapReduce permettant de produire les documents
   de type B à partir des documents de type A.


\begin{solution}
   Voici la modélisation en Javascript. Il faut, dans le Map, deux boucles imbriquées pour 
      accéder aux plateformes dont le nom constitue la clé de regroupement dans la phase de Reduce. 

    \begin{small}
        \begin{minted}{javascript}
      function mapAuteurVersPlateforme (doc) {
        // Première boucle sur les articles
        for article in doc.articles {
          // Second boucle sur les plateformes
          for (plateforme in article.plateforme) {
            // La clé est donc le nom de la platforme. On l'accompagne
            // d'un document contenant l'auteur, le titre et l'année
            emit (plateforme, {"auteur": doc.nom, "titre": article.titre, 
                                "année": article.année})
          }
         } 
       }
    \end{minted}
    \end{small}
    
    Pour le Reduce: on reçoit, pour chaque plateforme, la liste des articles avec leurs auteurs. 
    Il suffit de regrouper les articles par titre pour obtenir leurs auteurs.
    
        \begin{small}
    \begin{minted}{javascript}
      function reduceAuteurVersPlateforme (plateforme, [listeArticles]) {
         // On va s'épargner le petit algo de regroupement des auteurs sur le titre
         les_articles = groupement_auteurs_par_titre (listeArticles)
         // On a tout ce qu'il faut
         return ({"nom": platforme, "articles": les_articles} )
       }
    \end{minted}
    \end{small}
    
\end{solution}


   \textbf{Important}\,: pour spécifier les calculs MapReduce, donner les fonctions de Map et de Reduce, en javascript,
en pseudo-code, au pire en langage naturel en étant le plus précis possible. 
   

  \item On suppose que le contenu de la base  est accessible à un processeur
    Pig latin sous forme de trois collections. La première collection est celle
    des articles avec des documents du type\,:
    \begin{small}
     \begin{minted}{json}
{"ref" : "ju21", "titre": "D. Trump propose de racheter la Corse", "année": 2028}
    \end{minted}
    \end{small}
    La collection des auteurs. Exemples\,:
    \begin{small}
        \begin{minted}{json}
{"nom" : "P. Rigaux", "refArticle": "ju21"}
{"nom" : "P. Rigaux", "refArticle": "ju798"}
    \end{minted}
    \end{small}
    Et enfin la collection des diffusions. Exemples\,:
    \begin{small}
     \begin{minted}{json}
{"plateforme" : "Béru", "refArticle": "ju21"}
{"plateforme" : "Vautour", "refArticle": "ju798"}
    \end{minted}
    \end{small}
     
    On veut obtenir, à partir 
    de l'échantillon de base de données donné en début d'énoncé, un document
    représentant les informations sur un article\,:

    
    \begin{small}
    \begin{minted}{json}
{"ref": "ju21",
 "titre" : "D. Trump propose de racheter la Corse",
 "année": 2028,
 "auteurs": [
            ],
 "plateformes": [
         ]
}
    \end{minted}
    \end{small}
    
    
    \begin{enumerate}
      \item Donnez le contenu JSON des champs \texttt{auteurs} et \texttt{plateformes}
          d'après la base relationnelle de l'énoncé.
          
     \item Quelle chaîne de traitement permet de passer des trois collections
         à la représentation sous forme de document ci-dessus? Exprimez 
         la chaîne de traitement avec des opérateurs (pas besoin d'un syntaxe complète,
         un texte précis ou un schéma commenté
         fait très bien l'affaire).
      
\begin{solution}
L'opérateur principal à mettre en œuvre est le \texttt{cogroup}. Le premier
(mais l'ordre n'est pas important) groupe les articles et les auteurs.

\begin{verbatim}
articleEtAuteur = cogroup Articles by ref, Auteurs by refArticle;
\end{verbatim} 

Puis on regroupe la collection obtenue avec la collection des diffusions.


\begin{verbatim}
articlesComplet = cogroup articleEtAuteur by ref, Diffusions by refArticle;
\end{verbatim} 
\end{solution}     

\item À votre avis combien faut-il d'opérations MapReduce successives pour effectuer ce calcul\,?

\begin{solution}
Clairement, il en faut deux.
\end{solution}
    \end{enumerate}
\end{enumerate}

\begin{comment}
\section*{Deuxième partie : flux (4 pts)}

Chaque fois que quelqu'un lit un article sur une plateforme, un 
message est produit. Le format de ce message est le suivant\,:

  \begin{small}
    \begin{verbatim}
      {"refArticle": "ju9", "contribution": 0.005}
    \end{verbatim}
    \end{small}

La contribution est le montant (en Euros) généré par cette lecture. Ces messages
arrivent en  un flux nommé \texttt{lectures}. 
On veut écrire un traitement de flux qui cumule le montant des contributions par article. 

\begin{enumerate}
\item En vous inspirant des chaînes de traitement sur les flux vues dans le dernier chapitre
du cours, indiquez comment on maintient, pour chaque article, le cumul des contributions.

\begin{solution}
Voici la chaîne de traitement en Scala, décalquée à peu de choses près du compteur de mots.
\begin{small}
\begin{verbatim}
case class Lecture(refArticle: String, contribution: float)
val w = lectures.keyBy("refArticle") // Regroupement par article
        .sum("contribution")         // Cumul à chaque nouvelle lecture
        .print()              // Optionnel: affiche l'évolution continue du cumul
senv.execute("Cumul des contributions")
\end{verbatim}
\end{small}
\end{solution}

\item On veut rémunérer les auteurs d'articles toutes les 6 heures. Quelle est 
   la technique à utiliser pour obtenir le montant\,?
   
\begin{solution}
Il faut évidemment faire un fenêtrage. On maintient un accumulateur \texttt{accum}
dans lequel on ajoute les lectures reçues sur le flux, et ce pendant la durée de la fenêtre.
\begin{small}
\begin{verbatim}
case class Lecture(refArticle: String, contribution: float)
val w = lectures.keyBy("refArticle") // Regroupement par article
               .window(6*60*60) // Si on exprime en secondes
               .reduce( (accum, lecture) => (accum.ref, accum + lecture.contribution) }
               .print()
senv.execute("")
\end{verbatim}
\end{small}
\end{solution}

  \item Quelle est la parallélisation maximale permise par ce traitement\,?
\begin{solution}
Il y a un sous-flux par article, ce qui permet un large parallélisme à priori.
\end{solution}
\end{enumerate}
\end{comment}

\section*{Deuxième partie : recherche d'information (6 pts)}


Voici un échantillon du flux de nouvelles sur les éléctions régionales. 

\begin{enumerate}
 \item (d1) La Corse et l'Alsace sont deux régions pourvu d'une forte identité
     et d'une langue régionale. L'Alsace est plus calme...
        \item (d2) L'Alsace, la Corse, et la Normandie n'ont pas grand chose à voir avec la Bretagne.
        \item (d3) le nouveau président de la région Bretagne
                a prononcé son discours d'investiture.
           \item (d4) La Corse et la Normandie montrent la progression du parti 
                 des chaussettes vertes\,: analyse détaillée pour la Normandie.
\end{enumerate}

On veut indexer les nouvelles reçues de manière à pouvoir les rechercher
en fonction d'une région. Le vocabulaire auquel on se restreint est donc celui
des noms de pays (Alsace, Corse, Bretagne, Normandie).

\smallskip

\begin{enumerate}
   \item Donnez une matrice d'incidence contenant les tf pour les quatre régions,
      et un tableau donnant les idf (sans appliquer le $\log$). Mettez les noms de région
      en ligne, et les documents en colonne.

\begin{solution}
\begin{center}
\begin{tabular}{l|c|c|c|c}
            & \textbf{d1} & \textbf{d2} & \textbf{d3} &\textbf{d4} \\ \hline  
   Alsace  (2)   & 2          & 1           & 0           & 0        \\
   Corse    (4/3)  & 1           & 1           & 0          & 1         \\
   Bretagne    (2)  & 0           & 1          & 1           & 0         \\
   Normandie    (2)  & 0           & 1           & 0          & 2         \\ \hline
   \end{tabular}
\end{center}
\end{solution}

  \item Donner les résultats classés par similarité cosinus basée sur les tf 
  (on ignore l'idf) pour les requêtes suivantes. Expliquez brièvement le classement.
  
  \begin{itemize}
    \item Normandie;
%    \item Alsace et Corse ;
    \item Corse et Bretagne.
  \end{itemize}
\begin{solution}
Les normes
\begin{itemize}
  \item $||d_1|| = \sqrt{4+ 1} = \sqrt{5}$; $||d_2|| = \sqrt{1+1+1+1} = 2$; 
     $||d_3|| = \sqrt{1}$; 
    $||d_4|| = \sqrt{1+4}=\sqrt{5}$.
\end{itemize}

Les cosinus (requête non normalisée).

\begin{itemize}
  \item Normandie: d2 : $\frac{1}{2}=0,5$;  
       d4 : $\frac{2}{\sqrt{5}} \simeq 0,89$ 
       
       d4 est premier car il mentionne deux fois la Normandie.
  
  \item Alsace et Corse 
  
    d1 : $\frac{2+1}{\sqrt{5}}$; 
    d2 : $\frac{1+1}{2}$;
     d4 : $\frac{1}{\sqrt{5}}$; 
    
    d1 est premier comme on pouvait s'y attendre: il parle exclusivement
    de la Corse et de l'Alsace. Viennent ensuite d2 puis d4.
    
  \item Corse et Bretagne.
  
    d1 : $\frac{1}{\sqrt{5}}$; 
    d2 : $\frac{2}{2}$;
    d3 :  $\frac{1}{1}$  
    d4 :  $\frac{1}{\sqrt{5}}$  
    
  d2 et d3 arrivent à égalité. Intuitivement, il parlent tous les deux "à moitié"
   de Corse et de Bretagne.  d1 et d4 parlent de Corse \emph{ou} de Bretagne et aussi des autres
    pays.
 \end{itemize}

\end{solution}  
    
    \item Reprenons la dernière requête, "Corse et Bretagne" et les documents d2 et d3.
       Qu'est-ce que la prise en compte de l'idf changerait au classement?
    
\begin{solution}
La Bretagne a un idf plus élevé: d3 serait donc considéré come plus pertinent, 
car plus spécifique au besoin exprimé par la requête.
\end{solution}

  \end{enumerate}


  \section{Troisième partie\,: comprendre Cassandra (4 pts)}

  On veut proposer un schéma Cassandra pour explorer cette base. 
  Rappelons d'abord ce qu'est une clé composite en Cassandra\,: 
  \textit{A compound primary key consists of the partition key and the clustering key. The partition key determines 
  which node stores the data. Rows 
  for a partition key are stored in order based on the clustering key.} 

  On crée la table Article avec la commande suivante\,:
\begin{minted}{sql}
  create table Article (id_plateforme text,
                        id_article text,
                         auteur frozen<Personne>,
                         contenu text,
                         primary key (id_plateforme, id_article));
  \end{minted}

Commentez chacune des affirmations suivantes, indiquez si elles
sont vraies ou fausses et pourquoi.

  \begin{enumerate}
	\item tous les articles d'une plateforme sont sur le même serveur
    \item tous les articles d'un auteur sont sur le même serveur
	\item il n'y a pas de duplication d'un article, même s'il
	    a été publié par plus d'une plateforme
    \item Accéder aux articles d'une plateforme peut s'effectuer
	    par un parcours séquentiel sur un seul serveur
  \end{enumerate}  
  

  \begin{solution}
  \begin{enumerate}
    \item Oui puisque les articles d'une plateforme ont la même clé de partitionnement.
     \item Non, il n'y a aucun regroupement des articles d'un même auteur
     \item Si, il y a duplication d'un article autant de fois qu'il a été publié
     \item Oui !
  \end{enumerate} 
\end{solution}

\section*{Quatrième partie\,: systèmes distribués (3 pts)}


\begin{enumerate}

  \item Que signifie, pour  une structure de partitionnement, être \emph{dynamique}. Avez-vous
     étudié un système où le partitionnement n'était pas dynamique\,?
   
	 \item Quel est le principe de la reprise sur panne dans MapReduce?\,

	    \item Quel est le principe de la reprise sur panne dans Spark?\,
   
   \item Un groupe de serveurs en mode maître-esclave comprend 5 n{\oe}uds. Suite à une fragmentation
du réseau, le maître se retrouve  avec un seul esclave.  Expliquez ce qui se passe alors.
   
\end{enumerate}

\end{document}
